\documentclass{article}

\usepackage{amsmath,amssymb}
\usepackage{xurl}
\usepackage[colorlinks=true,linkcolor=blue,citecolor=blue,urlcolor=blue]{hyperref}
\setlength{\emergencystretch}{2em}

\input{command}

\title{The polynomial definition: functions with a representative,
not function types}
\author{}
\date{\today}

\begin{document}
\maketitle

\section{The paper's definition}

In \cite{Ji2020LowIndividualDegree}, an $m$-variate polynomial over
$\F_q$ is defined as a \emph{function}
$g : \F_q^m \to \F_q$ of the form
\[
  g(x_1,\dots,x_m) = \sum_{i_1,\dots,i_m}
    c_{i_1,\dots,i_m} \cdot x_1^{i_1}\cdots x_m^{i_m},
\]
with coefficients $c_{i_1,\dots,i_m} \in
\F_q$.\footnote{See \path{references/ldt-paper/introduction.tex},
line~3.}
Individual degree $d$ means $i_1,\dots,i_m \le d$ for every nonzero
coefficient. The set
\[
  \polyfunc{m}{q}{d}
  \subset  \{\, \F_q^m \to \F_q \}
\]
is then the set of functions that admit such a
representation.\footnote{See
\path{references/ldt-paper/preliminaries.tex}, lines~92--93.}
This is the standard convention in theoretical computer science:
``polynomial'' means a function that can be written as an algebraic
polynomial.

\section{The formal definition}

The formalization indexes measurement outcomes by algebraic
polynomials of low individual degree.\footnote{The outcome type is the
structure \leanid{MIPStarRE.LDT.Polynomial} in
\path{MIPStarRE/LDT/Basic/LowDegreePolynomial.lean}; its evaluation is
\leanid{MIPStarRE.LDT.Polynomial.toFun}.}  An outcome is a pair
consisting of
\[
  p \in \F_q[x_1,\dots,x_m]
  \quad\text{and a proof of}\quad
  \forall i,\ \deg_{x_i}(p) \le d,
\]
and its associated function $\F_q^m \to \F_q$ is the evaluation of
$p$.  Here $p$ ranges over the algebraic polynomial ring over the chosen
finite field,\footnote{This is the abbreviation
\leanid{PolynomialModel}, which expands to
\leanid{MvPolynomial (Fin m) (Scalar params)} in
\path{MIPStarRE/LDT/Basic/ParametersBase.lean}.}
and evaluation is coefficient-wise.  A function lies in
$\polyfunc{m}{q}{d}$ exactly when it is the evaluation of such an outcome.
When $d<q$, evaluation is injective on polynomials of individual degree at
most $d$, so the outcome set is in bijection with $\polyfunc{m}{q}{d}$.
When $d\ge q$ the error parameter of the main theorem is at least one, and
the theorem holds for trivial reasons.

The earlier blueprint used a function-type representation for
$\polyfunc{m}{q}{d}$, treating low-degree polynomials as a subtype
of the function space. That representation has been corrected to
algebraic representatives, which matches the paper's definition in the
regime $d<q$.\footnote{The blueprint correction is visible in
\path{blueprint/src/chapter/ch02_test.tex} and
\path{blueprint/src/chapter/ch03_preliminaries.tex}, where
$\polyfunc{m}{q}{d}$ is now the set of functions representable as an
algebraic polynomial of individual degree $\le d$.}

\section{Why algebraic polynomials?}

The paper's proofs (and the formalization) require algebraic
polynomials for the Schwartz--Zippel lemma.\footnote{Labelled
\leanid{lem:schwartz-zippel-total-degree} in the paper, and available
in Mathlib as \leanid{MvPolynomial.schwartzZippel}.}
The Mathlib statement operates on algebraic polynomials, not on
arbitrary functions. Indexing outcomes by algebraic representatives makes the
representative available wherever Schwartz--Zippel is applied.

The measurement families $G$, $A$, $B$, $L$ in the formalization
carry outcomes indexed by \emph{polynomials}, and the evaluation
maps $g \mapsto G_{g}$ respect the algebraic polynomial structure.
This design closely follows the paper, where measurement outcomes
are indexed by $g \in \polyfunc{m}{q}{d}$ and the Schwartz--Zippel
step uses algebraic properties of the polynomials.

\section{Conclusion}

There is no divergence in the relevant regime. The paper defines
polynomials as functions with a polynomial representative; the
formalization indexes outcomes by algebraic representatives of individual
degree at most $d$. For $d<q$ the two index sets are in bijection, and for
$d\ge q$ the main theorem is trivial.

\bibliographystyle{alpha}
\bibliography{references}

\end{document}